\documentclass[a4paper,10pt]{article}
\usepackage[utf8]{inputenc}
\usepackage[spanish]{babel}
\usepackage[T1]{fontenc}
\usepackage{geometry,enumitem,titlesec,parskip,microtype}
\usepackage[hidelinks]{hyperref}
\geometry{left=0.5in,right=0.5in,top=0.4in,bottom=0.4in}
\titlespacing*{\section}{0pt}{7pt}{3pt}
\titleformat{\section}{\normalfont\large\bfseries}{}{0pt}{}
\setlist[itemize]{leftmargin=12pt,topsep=2pt,parsep=0pt,itemsep=1pt}
\setlist[itemize,2]{label=--,leftmargin=10pt,topsep=1pt,itemsep=1pt}
\setlength{\parskip}{2pt}
\pagestyle{empty}
\begin{document}
\begin{center}
{\LARGE\bfseries Benjamin Schindler}\\[3pt]
\textbf{Líder de Ingeniería de IA} $\cdot$ Magíster en Ciencia de Datos\\[4pt]
+56 9 5627 9434 $\cdot$ \href{mailto:benjamin.schindlerv@gmail.com}{benjamin.schindlerv@gmail.com} $\cdot$ Santiago, Chile\\[2pt]
\href{https://benjaminschindler.vercel.app}{benjaminschindler.vercel.app} $\cdot$ \href{https://github.com/BenjaSchindler}{github.com/BenjaSchindler} $\cdot$ \href{https://www.linkedin.com/in/benjamin-schindler-92881a2b2/}{LinkedIn}
\end{center}
\section*{Perfil}
Soy ingeniero de IA y CTO en Doctor911, donde lidero un equipo de producto de tres personas. Desarrollo asistentes de IA, sistemas de búsqueda con RAG y herramientas para automatizar tareas del negocio.
\section*{Experiencia}
\textbf{Doctor911} \hfill Mar 2025 -- Presente\\
\textit{Chief Technology Officer (CTO)} \hfill \textit{Ene 2026 -- Presente}
\begin{itemize}
\item Lidero un equipo de producto de tres personas: dos ingenieros y un diseñador UX/UI.
\item Los productos de IA que desarrollé contribuyeron al crecimiento de 10$\times$ del tráfico web y a la rentabilidad.
\item Desarrollé triaje por voz con Gemini Live y lectura de exámenes con Cloud Vision OCR.
\item Diseñé la búsqueda RAG en Vertex AI y un servidor MCP para herramientas internas.
\item Lidero la estrategia de IA de Doctor911, seleccionada para ChileMass Emprende 2026.
\end{itemize}
\textit{Ingeniero de IA} \hfill \textit{Mar 2025 -- Ene 2026}
\begin{itemize}
\item Construí cuatro agentes de WhatsApp y dos agentes web con LangGraph y FastAPI.
\item Integré recomendaciones de exámenes, Meta Flows, catálogos y pagos con Transbank.
\item Combiné agentes con reglas de negocio y revisión médica para las aprobaciones clínicas.
\item Desplegué la plataforma en GCP Cloud Run con monitoreo de producción.
\end{itemize}
\textbf{WiseConn Latam} \hfill Sep 2024 -- Mar 2025\\
\textit{Practicante de Ciencia de Datos}
\begin{itemize}
\item Desarrollé predicciones de riego con CNN y clasificadores para instalaciones solares.
\item Desplegué modelos en SageMaker y S3 con validación y seguimiento de métricas.
\item Lideré una iniciativa de gobierno de datos entre departamentos.
\end{itemize}
\textbf{Unitti} \hfill Jul 2022 -- Jul 2024\\
\textit{Ingeniero de IA Jr.}
\begin{itemize}
\item Automaticé la extracción y conciliación de unas 800 facturas PDF al mes.
\item 85\% de las facturas sin intervención manual; procesamiento de diez a menos de un minuto y cierre mensual de ocho a tres días.
\item Desarrollé una aplicación de lenguaje natural a SQL y ajusté las consultas con comentarios de usuarios.
\item Desarrollé APIs en Python sobre PostgreSQL con pruebas automatizadas.
\end{itemize}
\section*{Proyectos}
\textbf{MiAutoCheck} (cliente): Una API que combina fotografías del vehículo e información de mercado en un informe PDF de tasación. Inspección visual, cinco agentes de investigación y un supervisor.

\textbf{EPE} (cliente): Funciones de IA para bienestar laboral: análisis de emociones y ejercicios personalizados, con detección de crisis y ocultamiento de datos personales. Prompts A/B, LLM-as-judge y Langfuse.

\textbf{Operaciones de Doctor911}: agente para coordinar al equipo y automatizar tareas recurrentes.
\section*{Educación}
\textbf{Magíster en Ciencia de Datos} $\cdot$ U. Adolfo Ibáñez \hfill 2024 -- 2026\\
\textbf{Ingeniería Civil Informática} $\cdot$ U. Adolfo Ibáñez \hfill 2020 -- Jul 2025\\
\textbf{Ingeniería Civil Industrial} $\cdot$ U. Adolfo Ibáñez \hfill 2020 -- Jul 2025\\
Magíster: 6,23/7,0. Tesis con Distinción Máxima: +2,25 pp macro-F1 sobre SMOTE (3.675 configuraciones).
\section*{Herramientas e idiomas}
\textbf{IA:} LangGraph, LangChain, RAG, MCP, PyTorch, XGBoost, OpenAI, Anthropic, Gemini Live, Langfuse.\\
\textbf{Desarrollo:} Python, TypeScript, SQL, FastAPI, Next.js, PostgreSQL, Docker, CI/CD.\\
\textbf{Cloud:} GCP (Cloud Run, Vertex AI, Cloud Vision), AWS (SageMaker, S3), Azure.\\
\textbf{Idiomas:} Español (nativo), Inglés (avanzado), Alemán (básico).
\end{document}
